\scanpage{241}\mspage{96}
2) Un cas particulièrement intéressant \inserted{de 7.4}
est celui où on suppose que les $\,{}_{y\backslash}X$ $(y\in\operatorname{Ob}Y)$
sont \deleted{\unsure{toutes}} sommes directes de catégories
ayant chacune un \uline{objet initial}. Alors
on voit que cette condition est stable par
les opérations de 7.4. a), et \deleted{\unsure{que}}\unsure{par}
conséquent, les colocalisés \inserted{$\,{}_{z\backslash}(X_{/x})$} des morphismes
localisés en haut $X_{/x}\unsure{-|\to}Y_{/y}$ sont des
catégories à objet initial. Ceci \deleted{\unsure{reste}
\unsure{satisfait}} implique alors qu'il en est de
$=$ des colocalisés $\,{}_{z\backslash}U$ $(z\in\operatorname{Ob}V)$ dans
$g:U\unsure{-|\to}V$. Il s'ensuit que $\unsure{S|g}$ est
$W$-coasphérique pour \uline{tt} \deleted{\unsure{loc}}\unsure{localiseur} \unsure{''}
fondamental $W$ (sans avoir à se servir
de l'axiome des limites).

La condition envisagée sur $f:X\unsure{-|\to}Y$
est satisfaite notamment dans les
deux cas suivants (qui me paraissent
de loin les plus importants) : a) $f$ est
étale et b) $X$ est discrète.
% archive sheet number: 241
